% !PRINTED 311
\weberpage{311}
%% [page furniture, not body text -- commented out 2026-08-10] Page furniture (running
%% head, folio, signature) is excluded from this edition by standing policy.
%% \noindent \S.\,81. \hfill Tripelsysteme.

%<<CONT
für drei und nur für drei Systeme der Unbekannten $\xi_1, \xi_2, \xi_3$ befriedigt werden, und dass diese drei ein Tripel bilden, so dass das Tripel durch die drei Congruenzen
\weq{5}{\Sigma \alpha\xi \equiv 0, \quad \Sigma \alpha\eta \equiv 0, \quad \Sigma \alpha\zeta \equiv 0 \quad (\mathrm{mod}\ 2)}
bestimmt ist. Man erkennt dies ohne weitläufige allgemeine Betrachtungen, wenn man die Fälle, in denen nur ein $\alpha$ oder zwei oder alle drei $\alpha \equiv 1$ (mod 2) sind, einzeln betrachtet.

Da es sieben verschiedene Zahlensysteme $\alpha_1, \alpha_2, \alpha_3$ giebt, so ist durch die Relationen (5) aus jedem solchen Systeme der $\alpha$ ein Tripel völlig bestimmt, und wir können die Function $v$ eindeutig durch $v_{\alpha_1,\alpha_2,\alpha_3}$ oder kürzer durch $v_\alpha$ bezeichnen:
\weq{6}{v_\alpha = \frac{1}{\Psi_\xi \Psi_\eta \Psi_\zeta}.}

\gesperrt{3. Die Functionen $v_\alpha$, als Functionen der $X$ aufgefasst, gestatten die Permutationen der Gruppe $Q$.}

Wenn wir auf die $z$ eine Substitution $\sigma$ aus $S$ anwenden, so erleiden nach 2.\ (\S.\,80) die $\xi, \eta, \zeta$ gleichzeitig die Substitution $\varrho^{-1}$; daraus aber ergiebt sich nach (5), dass die $\alpha$ die Substitution $\sigma$ erleiden, und wir haben also:

\gesperrt{4. Die Anwendung der Substitution $\sigma$ aus $S$ auf die Indices $z$ von $X$ hat die Anwendung derselben Substitution auf die Indices $\alpha$ von $v_\alpha$ zur Folge.}

Setzen wir in den Summen (5), in denen $(\xi), (\eta), (\zeta)$ das zu $(\alpha)$ gehörige Tripel ist, für $(\alpha)$ ein anderes System $(\beta)$, so ist wegen (3)
\[
\Sigma \beta\xi + \Sigma \beta\eta + \Sigma \beta\zeta \equiv 0.
\]

Ist nun $\beta$ von $\alpha$ verschieden, so können diese drei Summen nicht alle drei congruent mit 0 sein, weil sonst zwei verschiedene Systeme $(\alpha), (\beta)$ zu demselben Tripel führen würden; \gesperrt{folglich müssen von diesen Summen zwei ungerade und eine gerade sein.} Wir haben daher den Satz:

\gesperrt{5. Durchläuft $(\xi)$ das zu $(\alpha)$ gehörige Tripel, und ist $(\beta)$ ein von $(\alpha)$ verschiedenes System, so sind unter den Summen $\Sigma \beta\xi$ zwei ungerade und eine gerade.}
